% ============================================================================
\clearpage
\section*{v4.9.0-r35 Source-Only: Inner and Outer Complex Unit in the Colour Assignment}
\addcontentsline{toc}{section}{v4.9.0-r35 Source-Only: Inner and Outer Complex Unit in the Colour Assignment}
\label{sec:v490-r35-inner-outer-unit}
% ============================================================================

\paragraph{Purpose.}
Layer v4.9.0-r34 refuted the literal colour and \(\nu_R\) assignment of Theorem~\ref{thm:octonion-SU3} and
Theorem~\ref{thm:three-generations} (r34.1(c)).  At the author's suggestion that this may again be a fixed-frame effect, this layer
records the pre-registered test H20 (note \texttt{llm/Vorab\_Notiz\_H20\_Frame\_Zuordnung.md}, commit \texttt{7d95feb}, fixed before
computing).  It separates two complex units: the \emph{inner} unit \(J=L_{e_7}\) on \(\OO\) itself, and the \emph{outer} unit \(i\),
the \(\mathbb C\) factor of the Dixon storage, on which the charge operator \(Q\) of r34 is built.  Earlier layers are not rewritten.

\noindent\textbf{Source-only status.}
This v4.9.0-r35 layer is source-only.  It does not build a PDF, does not assign a new DOI,
and does not authorize a public release.  The v4.9.0 Version DOI
\href{https://doi.org/10.5281/zenodo.23534954}{10.5281/zenodo.23534954} and the Concept DOI
\href{https://doi.org/10.5281/zenodo.19210728}{10.5281/zenodo.19210728} are unchanged by it.

\subsection*{v4.9.0-r35 results}
\label{subsec:v490-r35-results}

Conventions as in v4.9.0-r34.  Put \(P_\pm=\tfrac12(1\mp i\,J)\colon\OO\to\mathbb C\otimes\OO\), and let \(R_{e_7}\) be right
multiplication by \(e_7\).

\paragraph{Theoremlet v4.9.0-r35.1.}
\begin{enumerate}[leftmargin=2em]
\item[(a)] \emph{Inner frame.}  \(\operatorname{Stab}_{G_2}(e_7)\) is \(J\)-linear.  \((\OO,J)\cong\mathbb C^4\) splits into the trivial
line \(\operatorname{span}_{\mathbb R}(1,e_7)\) and the irreducible complex triplet \(e_7^\perp\), with complex basis \(e_4,e_5,e_6\);
\(Je_4=-e_3\), \(Je_5=-e_2\), \(Je_6=e_1\).
\item[(b)] \(\operatorname{span}(e_1,e_2,e_3)=J\operatorname{span}(e_4,e_5,e_6)\) is neither \(J\)-invariant nor invariant under
\(\operatorname{Stab}_{G_2}(e_7)\) nor under \(Q\).
\item[(c)] \(P_\pm\) is \(\operatorname{Stab}_{G_2}(e_7)\)-equivariant and injective, and \(P_\pm J=\pm i\,P_\pm\).  \(P_-\) maps
\(e_7^\perp\) onto the triplet \(Q=\tfrac23\) and \(\operatorname{span}(1,e_7)\), in particular \(e_7\), onto the singlet \(Q=0\);
\(P_+\) maps them onto \(Q=\tfrac13\) and \(Q=1\).  \(T\) exchanges \(P_+\) and \(P_-\).
\item[(d)] \(Q=\tfrac12-\tfrac{i}{6}\,L_{e_7}-\tfrac{i}{3}\,R_{e_7}\).
\end{enumerate}

\paragraph{Proof sketch.}
(a)--(c) by explicit matrices (script \texttt{llm/k76\_test\_h20.py}, deviations at most \(2\cdot10^{-15}\)); irreducibility of complex
type from the real commutant on \(e_7^\perp\), which is two-dimensional.  In (c), \(P_\pm\) commutes with every \(J\)-linear map, and
\(P_\pm J=\tfrac12(J\pm i)=\pm i\,P_\pm\); \(\bar P_\pm=P_\mp\) because \(J\) is real.  (d) by a least-squares fit with residual
\(2\cdot10^{-16}\).  \(\square\)

\paragraph{Verdict and reading.}
H20a, H20c and H20d are confirmed; H20b is refuted.
\begin{enumerate}[leftmargin=2em]
\item The triplet basis \(e_4,e_5,e_6\) and \(\nu_R\leftrightarrow e_7\) are correct in the inner frame, and
Theorem~\ref{thm:octonion-SU3} holds there.  H19c in r34 took the real vectors as states of \(\mathbb C\otimes\OO\), which mixes up
the two frames.  The translation to the outer frame is \(P_\pm\), and the sign is a frame choice.  With \(P_-\) the Companion's
assignment agrees with Furey's reading: \(e_7\) goes to the singlet \(Q=0\) (\(\nu\)) and \(e_4,e_5,e_6\) to the triplet
\(Q=\tfrac23\) (\(u\)).  With \(P_+\) one obtains the conjugate picture.  This is r34.1(d) again: \(T\) exchanges \(P_\pm\), hence
\(Q\mapsto1-Q\).
\item The sentence in the proof of Theorem~\ref{thm:three-generations} that the anti-triplet \((e_1,e_2,e_3)\) carries the
conjugate colour representation is a genuine error.  It arises from mixing the frames: in the inner frame there is one triplet,
and \(e_1,e_2,e_3\) are its \(J\)-images; a conjugate triplet exists only with the outer unit, as the image of \(P_+\) against that
of \(P_-\), and not on \(\operatorname{span}(e_1,e_2,e_3)\).
\item The charge needs both units, and both left and right multiplication (r35.1(d)).  The complex scalar \(it\) of
Theorem~\ref{thm:SM} alone generates a global phase that is the same on all four charge sectors; it is not electromagnetism.
\item The double use of \(e_4,e_5,e_6\) for colour and generation (v4.8.0-r13) is not a frame effect and stays.
\item The \emph{reading} of r34.1(c) is corrected: the statement stands, but the Companion does not put colour and \(\nu_R\) on
wrong subspaces.  It states them in the inner frame without naming the translation \(P_\pm\).
\end{enumerate}
Proposed addition to Konvention~F: statements about colour and charge name their complex unit (inner \(J=L_{e_7}\), without
charge; outer \(i\), with \(Q\)), and for the translation the sign \(P_+\) or \(P_-\).

\subsection*{v4.9.0-r35 status in the three frames}
\label{subsec:v490-r35-frames}

\begin{center}
\small
\begin{tabular}{@{}>{\raggedright\arraybackslash}p{0.30\textwidth}>{\raggedright\arraybackslash}p{0.62\textwidth}@{}}
\toprule
result & M / T / P\\
\midrule
r35.1(a), (b) & inner frame only; no use of the \(\mathbb C\) factor, hence the same in M, T and P\\
r35.1(c), (d) & stated in M; \(T\) exchanges \(P_+\) and \(P_-\), so the sign of the translation is the M/T choice; P not evaluated\\
\bottomrule
\end{tabular}
\end{center}

\subsection*{v4.9.0-r35 register entries}
\label{subsec:v490-r35-register}

\begin{center}
\small
\begin{tabular}{@{}>{\raggedright\arraybackslash}p{0.34\textwidth}>{\raggedright\arraybackslash}p{0.58\textwidth}@{}}
\toprule
statement & status after v4.9.0-r35\\
\midrule
Theorem~\ref{thm:octonion-SU3} & holds in the inner frame (r35.1(a))\\
Theorem~\ref{thm:three-generations}, triplet basis and \(\nu_R\leftrightarrow e_7\) & correct in the inner frame; with \(P_-\) consistent with Furey's charges (r35.1(c)); reading of r34.1(c) corrected\\
Theorem~\ref{thm:three-generations}, anti-triplet \((e_1,e_2,e_3)\) & refuted, frame mixing (r35.1(b)); candidate for a status note\\
Theorem~\ref{thm:SM}, \(U(1)\) line & the scalar \(it\) gives only a global phase; the charge needs both units (r35.1(d))\\
\bottomrule
\end{tabular}
\end{center}

\subsection*{v4.9.0-r35 ledger and guard update}
\label{subsec:v490-r35-ledger-guard}

\[
R_{\rm fr}^{v490\text{-}r35}=(1,0,1)
\]
for (colour assignment correct in the inner frame; anti-triplet sentence wrong; frame rule for colour and charge proposed).

The layer spends no guard:
\[
R_{\rm glob}^{(v490\text{-}r35,+)}=(0,1,1,1,1),
\qquad
\Gamma_{\rm guard}^{v490\text{-}r35}=(1,1,1,1,1),
\qquad
\beta_4=1/256\ \text{unspent}.
\]
It does not imply exact finite \(C5\) zero closure, exact all-readout terminality,
\(C6/T2\), a global photon-clock/RH theorem, or global all-readout closure.
